% Lazo de calculo del flujo de calor que alcanza el MDNBR objetivo.
%
% Fragmento para % Lazo de calculo del flujo de calor que alcanza el MDNBR objetivo.
%
% Fragmento para % Lazo de calculo del flujo de calor que alcanza el MDNBR objetivo.
%
% Fragmento para % Lazo de calculo del flujo de calor que alcanza el MDNBR objetivo.
%
% Fragmento para \input{Figures/dnb_dnn_loop}.
% Requiere en el preambulo del documento principal:
%   \usepackage{tikz}
%   \usetikzlibrary{arrows.meta}
%   \usepackage{amsmath,amssymb}
%
% Estructura: condiciones de operacion -> ramas paralelas de piso (q''_min) y
% techo (q''_max), cada una con establecer flujo, calcular caudal masico y
% predecir MDNBR con la DNN -> lazo de la secante, con el calculo del caudal
% masico y la prediccion de la DNN como pasos separados, hasta cumplir el
% criterio de convergencia.
%
% Los tres pasos que invocan a la red llevan un esquema en miniatura de la DNN
% (pic "dnnmini"), consistente con fig:mlp-arquitectura.

\begin{figure}[htbp]
  \centering
  \resizebox{\textwidth}{!}{
  \begin{tikzpicture}[
      caja/.style   = {rectangle, rounded corners=2pt, draw=black!70,
                       fill=black!4, line width=0.6pt, text width=6.0cm,
                       align=center, font=\footnotesize, minimum height=1.0cm,
                       inner sep=3pt},
      subcaja/.style= {rectangle, rounded corners=2pt, draw=black!70,
                       fill=white, line width=0.6pt, text width=4.6cm,
                       align=center, font=\footnotesize, minimum height=1.0cm,
                       inner sep=3pt},
      final/.style  = {rectangle, rounded corners=2pt, draw=black!80,
                       fill=black!8, line width=0.9pt, text width=2.6cm,
                       align=center, font=\footnotesize, minimum height=0.85cm,
                       inner sep=3pt},
      grupo/.style  = {rectangle, rounded corners=3pt, draw=black!45,
                       line width=0.6pt, dashed},
      titulo/.style = {font=\footnotesize\bfseries, inner sep=1pt},
      rombo/.style  = {draw=black!70, fill=black!4, line width=0.6pt},
      texto/.style  = {font=\footnotesize, align=center, inner sep=1pt},
      nombre/.style = {font=\footnotesize, inner sep=1pt},
      flecha/.style = {-{Latex[length=2mm,width=1.6mm]}, draw=black!70,
                       line width=0.6pt},
      linea/.style  = {draw=black!70, line width=0.6pt},
      %--- estilos del esquema en miniatura de la red
      nnnodo/.style = {circle, draw=black!65, fill=white, line width=0.3pt,
                       minimum size=1.4mm, inner sep=0pt},
      nnarco/.style = {draw=black!30, line width=0.25pt},
      %--- esquema en miniatura de la DNN (3 entradas, 4+4 ocultas, 2 salidas)
      dnnmini/.pic  = {
        \foreach \ya in {0.20,0,-0.20}
          \foreach \yb in {0.30,0.10,-0.10,-0.30}
            \draw[nnarco] (0,\ya) -- (0.4,\yb);
        \foreach \ya in {0.30,0.10,-0.10,-0.30}
          \foreach \yb in {0.30,0.10,-0.10,-0.30}
            \draw[nnarco] (0.4,\ya) -- (0.8,\yb);
        \foreach \ya in {0.30,0.10,-0.10,-0.30}
          \foreach \yb in {0.10,-0.10}
            \draw[nnarco] (0.8,\ya) -- (1.2,\yb);
        \foreach \y in {0.20,0,-0.20}
          \node[nnnodo] at (0,\y) {};
        \foreach \y in {0.30,0.10,-0.10,-0.30} {
          \node[nnnodo] at (0.4,\y) {};
          \node[nnnodo] at (0.8,\y) {};
        }
        \foreach \y in {0.10,-0.10}
          \node[nnnodo] at (1.2,\y) {};
      },
    ]
    \useasboundingbox (-6.25,-15.80) rectangle (8.61,1.0);

    %-------------------------------------------------- 0. condiciones
    \node[caja] (A) at (0,0)
      {0.~Definir las condiciones de operaci\'on\\
       y el objetivo $\mathrm{MDNBR}^{*}$};

    %-------------------------------------------------- marcos de las ramas
    \draw[grupo] (-5.7,-6.97) rectangle (-0.3,-1.35);
    \draw[grupo] ( 0.3,-6.97) rectangle ( 5.7,-1.35);

    \node[titulo] at (-3.0,-1.96) {1.~C\'alculo de piso};
    \node[titulo] at ( 3.0,-1.96) {2.~C\'alculo de techo};

    %-------------------------------------------------- rama de piso
    \node[subcaja] (P1) at (-3.0,-3.07)
      {1.a~~Establecer el flujo\\ de calor m\'inimo $q''_{\min}$};
    \node[subcaja] (P2) at (-3.0,-4.42)
      {1.b~~Calcular el caudal\\ m\'asico $G$ para $q''_{\min}$};
    \node[subcaja, minimum height=1.45cm] (P3) at (-3.0,-5.995) {};
    \node[texto, text width=2.5cm] at (-3.90,-5.995)
      {1.c~~Predecir $\mathrm{MDNBR}_{\min}$ con la DNN};
    \pic at (-2.12,-5.995) {dnnmini};

    %-------------------------------------------------- rama de techo
    \node[subcaja] (T1) at (3.0,-3.07)
      {2.a~~Establecer el flujo\\ de calor m\'aximo $q''_{\max}$};
    \node[subcaja] (T2) at (3.0,-4.42)
      {2.b~~Calcular el caudal\\ m\'asico $G$ para $q''_{\max}$};
    \node[subcaja, minimum height=1.45cm] (T3) at (3.0,-5.995) {};
    \node[texto, text width=2.5cm] at (2.10,-5.995)
      {2.c~~Predecir $\mathrm{MDNBR}_{\max}$ con la DNN};
    \pic at (3.88,-5.995) {dnnmini};

    %-------------------------------------------------- lazo de la secante
    \node[caja] (C) at (0,-8.40)
      {3.~M\'etodo de la secante:\\
       nuevo flujo de calor $q''_{k+1}$};

    \node[caja, minimum height=0.85cm] (D1) at (0,-9.675)
      {3.a~~Calcular el caudal m\'asico $G_{k+1}$};

    \node[caja, minimum height=1.25cm] (D2) at (0,-11.075) {};
    \node[texto, text width=3.6cm] at (-0.92,-11.075)
      {3.b~~Predecir $\mathrm{MDNBR}_{k+1}$ con la DNN};
    \pic at (1.45,-11.075) {dnnmini};

    %-------------------------------------------------- criterio de convergencia
    \draw[rombo] (0,-12.05) -- (4.2,-13.65) -- (0,-15.25) -- (-4.2,-13.65)
      -- cycle;
    \node[texto] at (0,-13.65)
      {$\begin{aligned}
         \bigl|\mathrm{MDNBR}_{k+1}-\mathrm{MDNBR}^{*}\bigr|
           &\le\varepsilon_{\mathrm{M}} \ \text{y} \\[-1pt]
         \bigl|G_{k+1}-G_{k}\bigr| &\le\varepsilon_{G} \ \text{y} \\[-1pt]
         \bigl|q''_{k+1}-q''_{k}\bigr| &\le\varepsilon_{q} \ ?
       \end{aligned}$};

    \node[final] (F) at (6.705,-13.65) {$q''_{\mathrm{PLC}}=q''_{k+1}$};

    %-------------------------------------------------- flujo: 0 -> ramas
    \draw[linea]  (A.south) -- (0,-0.95);
    \draw[linea]  (-3.0,-0.95) -- (3.0,-0.95);
    \draw[flecha] (-3.0,-0.95) -- (-3.0,-1.35);
    \draw[flecha] ( 3.0,-0.95) -- ( 3.0,-1.35);

    %-------------------------------------------------- flujo dentro de cada rama
    \draw[flecha] (P1.south) -- (P2.north);
    \draw[flecha] (P2.south) -- (P3.north);
    \draw[flecha] (T1.south) -- (T2.north);
    \draw[flecha] (T2.south) -- (T3.north);

    %-------------------------------------------------- ramas -> secante
    \draw[linea]  (-3.0,-6.97) -- (-3.0,-7.50) -- (0,-7.50);
    \draw[linea]  ( 3.0,-6.97) -- ( 3.0,-7.50) -- (0,-7.50);
    \draw[flecha] (0,-7.50) -- (C.north);

    %-------------------------------------------------- lazo
    \draw[flecha] (C.south)  -- (D1.north);
    \draw[flecha] (D1.south) -- (D2.north);
    \draw[flecha] (D2.south) -- (0,-12.05);

    \draw[flecha] (4.2,-13.65) -- (F.west);
    \node[nombre] at (4.75,-13.20) {s\'i};

    \draw[linea]  (-4.2,-13.65) -- (-5.7,-13.65) -- (-5.7,-8.40);
    \draw[flecha] (-5.7,-8.40) -- (C.west);
    \node[nombre, anchor=west] at (-5.3,-11.0) {no};
  \end{tikzpicture}}
  \caption{Lazo de c\'alculo del flujo de calor que alcanza el MDNBR objetivo.}
  \label{fig:dnb-dnn-loop}
\end{figure}
.
% Requiere en el preambulo del documento principal:
%   \usepackage{tikz}
%   \usetikzlibrary{arrows.meta}
%   \usepackage{amsmath,amssymb}
%
% Estructura: condiciones de operacion -> ramas paralelas de piso (q''_min) y
% techo (q''_max), cada una con establecer flujo, calcular caudal masico y
% predecir MDNBR con la DNN -> lazo de la secante, con el calculo del caudal
% masico y la prediccion de la DNN como pasos separados, hasta cumplir el
% criterio de convergencia.
%
% Los tres pasos que invocan a la red llevan un esquema en miniatura de la DNN
% (pic "dnnmini"), consistente con fig:mlp-arquitectura.

\begin{figure}[htbp]
  \centering
  \resizebox{\textwidth}{!}{
  \begin{tikzpicture}[
      caja/.style   = {rectangle, rounded corners=2pt, draw=black!70,
                       fill=black!4, line width=0.6pt, text width=6.0cm,
                       align=center, font=\footnotesize, minimum height=1.0cm,
                       inner sep=3pt},
      subcaja/.style= {rectangle, rounded corners=2pt, draw=black!70,
                       fill=white, line width=0.6pt, text width=4.6cm,
                       align=center, font=\footnotesize, minimum height=1.0cm,
                       inner sep=3pt},
      final/.style  = {rectangle, rounded corners=2pt, draw=black!80,
                       fill=black!8, line width=0.9pt, text width=2.6cm,
                       align=center, font=\footnotesize, minimum height=0.85cm,
                       inner sep=3pt},
      grupo/.style  = {rectangle, rounded corners=3pt, draw=black!45,
                       line width=0.6pt, dashed},
      titulo/.style = {font=\footnotesize\bfseries, inner sep=1pt},
      rombo/.style  = {draw=black!70, fill=black!4, line width=0.6pt},
      texto/.style  = {font=\footnotesize, align=center, inner sep=1pt},
      nombre/.style = {font=\footnotesize, inner sep=1pt},
      flecha/.style = {-{Latex[length=2mm,width=1.6mm]}, draw=black!70,
                       line width=0.6pt},
      linea/.style  = {draw=black!70, line width=0.6pt},
      %--- estilos del esquema en miniatura de la red
      nnnodo/.style = {circle, draw=black!65, fill=white, line width=0.3pt,
                       minimum size=1.4mm, inner sep=0pt},
      nnarco/.style = {draw=black!30, line width=0.25pt},
      %--- esquema en miniatura de la DNN (3 entradas, 4+4 ocultas, 2 salidas)
      dnnmini/.pic  = {
        \foreach \ya in {0.20,0,-0.20}
          \foreach \yb in {0.30,0.10,-0.10,-0.30}
            \draw[nnarco] (0,\ya) -- (0.4,\yb);
        \foreach \ya in {0.30,0.10,-0.10,-0.30}
          \foreach \yb in {0.30,0.10,-0.10,-0.30}
            \draw[nnarco] (0.4,\ya) -- (0.8,\yb);
        \foreach \ya in {0.30,0.10,-0.10,-0.30}
          \foreach \yb in {0.10,-0.10}
            \draw[nnarco] (0.8,\ya) -- (1.2,\yb);
        \foreach \y in {0.20,0,-0.20}
          \node[nnnodo] at (0,\y) {};
        \foreach \y in {0.30,0.10,-0.10,-0.30} {
          \node[nnnodo] at (0.4,\y) {};
          \node[nnnodo] at (0.8,\y) {};
        }
        \foreach \y in {0.10,-0.10}
          \node[nnnodo] at (1.2,\y) {};
      },
    ]
    \useasboundingbox (-6.25,-15.80) rectangle (8.61,1.0);

    %-------------------------------------------------- 0. condiciones
    \node[caja] (A) at (0,0)
      {0.~Definir las condiciones de operaci\'on\\
       y el objetivo $\mathrm{MDNBR}^{*}$};

    %-------------------------------------------------- marcos de las ramas
    \draw[grupo] (-5.7,-6.97) rectangle (-0.3,-1.35);
    \draw[grupo] ( 0.3,-6.97) rectangle ( 5.7,-1.35);

    \node[titulo] at (-3.0,-1.96) {1.~C\'alculo de piso};
    \node[titulo] at ( 3.0,-1.96) {2.~C\'alculo de techo};

    %-------------------------------------------------- rama de piso
    \node[subcaja] (P1) at (-3.0,-3.07)
      {1.a~~Establecer el flujo\\ de calor m\'inimo $q''_{\min}$};
    \node[subcaja] (P2) at (-3.0,-4.42)
      {1.b~~Calcular el caudal\\ m\'asico $G$ para $q''_{\min}$};
    \node[subcaja, minimum height=1.45cm] (P3) at (-3.0,-5.995) {};
    \node[texto, text width=2.5cm] at (-3.90,-5.995)
      {1.c~~Predecir $\mathrm{MDNBR}_{\min}$ con la DNN};
    \pic at (-2.12,-5.995) {dnnmini};

    %-------------------------------------------------- rama de techo
    \node[subcaja] (T1) at (3.0,-3.07)
      {2.a~~Establecer el flujo\\ de calor m\'aximo $q''_{\max}$};
    \node[subcaja] (T2) at (3.0,-4.42)
      {2.b~~Calcular el caudal\\ m\'asico $G$ para $q''_{\max}$};
    \node[subcaja, minimum height=1.45cm] (T3) at (3.0,-5.995) {};
    \node[texto, text width=2.5cm] at (2.10,-5.995)
      {2.c~~Predecir $\mathrm{MDNBR}_{\max}$ con la DNN};
    \pic at (3.88,-5.995) {dnnmini};

    %-------------------------------------------------- lazo de la secante
    \node[caja] (C) at (0,-8.40)
      {3.~M\'etodo de la secante:\\
       nuevo flujo de calor $q''_{k+1}$};

    \node[caja, minimum height=0.85cm] (D1) at (0,-9.675)
      {3.a~~Calcular el caudal m\'asico $G_{k+1}$};

    \node[caja, minimum height=1.25cm] (D2) at (0,-11.075) {};
    \node[texto, text width=3.6cm] at (-0.92,-11.075)
      {3.b~~Predecir $\mathrm{MDNBR}_{k+1}$ con la DNN};
    \pic at (1.45,-11.075) {dnnmini};

    %-------------------------------------------------- criterio de convergencia
    \draw[rombo] (0,-12.05) -- (4.2,-13.65) -- (0,-15.25) -- (-4.2,-13.65)
      -- cycle;
    \node[texto] at (0,-13.65)
      {$\begin{aligned}
         \bigl|\mathrm{MDNBR}_{k+1}-\mathrm{MDNBR}^{*}\bigr|
           &\le\varepsilon_{\mathrm{M}} \ \text{y} \\[-1pt]
         \bigl|G_{k+1}-G_{k}\bigr| &\le\varepsilon_{G} \ \text{y} \\[-1pt]
         \bigl|q''_{k+1}-q''_{k}\bigr| &\le\varepsilon_{q} \ ?
       \end{aligned}$};

    \node[final] (F) at (6.705,-13.65) {$q''_{\mathrm{PLC}}=q''_{k+1}$};

    %-------------------------------------------------- flujo: 0 -> ramas
    \draw[linea]  (A.south) -- (0,-0.95);
    \draw[linea]  (-3.0,-0.95) -- (3.0,-0.95);
    \draw[flecha] (-3.0,-0.95) -- (-3.0,-1.35);
    \draw[flecha] ( 3.0,-0.95) -- ( 3.0,-1.35);

    %-------------------------------------------------- flujo dentro de cada rama
    \draw[flecha] (P1.south) -- (P2.north);
    \draw[flecha] (P2.south) -- (P3.north);
    \draw[flecha] (T1.south) -- (T2.north);
    \draw[flecha] (T2.south) -- (T3.north);

    %-------------------------------------------------- ramas -> secante
    \draw[linea]  (-3.0,-6.97) -- (-3.0,-7.50) -- (0,-7.50);
    \draw[linea]  ( 3.0,-6.97) -- ( 3.0,-7.50) -- (0,-7.50);
    \draw[flecha] (0,-7.50) -- (C.north);

    %-------------------------------------------------- lazo
    \draw[flecha] (C.south)  -- (D1.north);
    \draw[flecha] (D1.south) -- (D2.north);
    \draw[flecha] (D2.south) -- (0,-12.05);

    \draw[flecha] (4.2,-13.65) -- (F.west);
    \node[nombre] at (4.75,-13.20) {s\'i};

    \draw[linea]  (-4.2,-13.65) -- (-5.7,-13.65) -- (-5.7,-8.40);
    \draw[flecha] (-5.7,-8.40) -- (C.west);
    \node[nombre, anchor=west] at (-5.3,-11.0) {no};
  \end{tikzpicture}}
  \caption{Lazo de c\'alculo del flujo de calor que alcanza el MDNBR objetivo.}
  \label{fig:dnb-dnn-loop}
\end{figure}
.
% Requiere en el preambulo del documento principal:
%   \usepackage{tikz}
%   \usetikzlibrary{arrows.meta}
%   \usepackage{amsmath,amssymb}
%
% Estructura: condiciones de operacion -> ramas paralelas de piso (q''_min) y
% techo (q''_max), cada una con establecer flujo, calcular caudal masico y
% predecir MDNBR con la DNN -> lazo de la secante, con el calculo del caudal
% masico y la prediccion de la DNN como pasos separados, hasta cumplir el
% criterio de convergencia.
%
% Los tres pasos que invocan a la red llevan un esquema en miniatura de la DNN
% (pic "dnnmini"), consistente con fig:mlp-arquitectura.

\begin{figure}[htbp]
  \centering
  \resizebox{\textwidth}{!}{
  \begin{tikzpicture}[
      caja/.style   = {rectangle, rounded corners=2pt, draw=black!70,
                       fill=black!4, line width=0.6pt, text width=6.0cm,
                       align=center, font=\footnotesize, minimum height=1.0cm,
                       inner sep=3pt},
      subcaja/.style= {rectangle, rounded corners=2pt, draw=black!70,
                       fill=white, line width=0.6pt, text width=4.6cm,
                       align=center, font=\footnotesize, minimum height=1.0cm,
                       inner sep=3pt},
      final/.style  = {rectangle, rounded corners=2pt, draw=black!80,
                       fill=black!8, line width=0.9pt, text width=2.6cm,
                       align=center, font=\footnotesize, minimum height=0.85cm,
                       inner sep=3pt},
      grupo/.style  = {rectangle, rounded corners=3pt, draw=black!45,
                       line width=0.6pt, dashed},
      titulo/.style = {font=\footnotesize\bfseries, inner sep=1pt},
      rombo/.style  = {draw=black!70, fill=black!4, line width=0.6pt},
      texto/.style  = {font=\footnotesize, align=center, inner sep=1pt},
      nombre/.style = {font=\footnotesize, inner sep=1pt},
      flecha/.style = {-{Latex[length=2mm,width=1.6mm]}, draw=black!70,
                       line width=0.6pt},
      linea/.style  = {draw=black!70, line width=0.6pt},
      %--- estilos del esquema en miniatura de la red
      nnnodo/.style = {circle, draw=black!65, fill=white, line width=0.3pt,
                       minimum size=1.4mm, inner sep=0pt},
      nnarco/.style = {draw=black!30, line width=0.25pt},
      %--- esquema en miniatura de la DNN (3 entradas, 4+4 ocultas, 2 salidas)
      dnnmini/.pic  = {
        \foreach \ya in {0.20,0,-0.20}
          \foreach \yb in {0.30,0.10,-0.10,-0.30}
            \draw[nnarco] (0,\ya) -- (0.4,\yb);
        \foreach \ya in {0.30,0.10,-0.10,-0.30}
          \foreach \yb in {0.30,0.10,-0.10,-0.30}
            \draw[nnarco] (0.4,\ya) -- (0.8,\yb);
        \foreach \ya in {0.30,0.10,-0.10,-0.30}
          \foreach \yb in {0.10,-0.10}
            \draw[nnarco] (0.8,\ya) -- (1.2,\yb);
        \foreach \y in {0.20,0,-0.20}
          \node[nnnodo] at (0,\y) {};
        \foreach \y in {0.30,0.10,-0.10,-0.30} {
          \node[nnnodo] at (0.4,\y) {};
          \node[nnnodo] at (0.8,\y) {};
        }
        \foreach \y in {0.10,-0.10}
          \node[nnnodo] at (1.2,\y) {};
      },
    ]
    \useasboundingbox (-6.25,-15.80) rectangle (8.61,1.0);

    %-------------------------------------------------- 0. condiciones
    \node[caja] (A) at (0,0)
      {0.~Definir las condiciones de operaci\'on\\
       y el objetivo $\mathrm{MDNBR}^{*}$};

    %-------------------------------------------------- marcos de las ramas
    \draw[grupo] (-5.7,-6.97) rectangle (-0.3,-1.35);
    \draw[grupo] ( 0.3,-6.97) rectangle ( 5.7,-1.35);

    \node[titulo] at (-3.0,-1.96) {1.~C\'alculo de piso};
    \node[titulo] at ( 3.0,-1.96) {2.~C\'alculo de techo};

    %-------------------------------------------------- rama de piso
    \node[subcaja] (P1) at (-3.0,-3.07)
      {1.a~~Establecer el flujo\\ de calor m\'inimo $q''_{\min}$};
    \node[subcaja] (P2) at (-3.0,-4.42)
      {1.b~~Calcular el caudal\\ m\'asico $G$ para $q''_{\min}$};
    \node[subcaja, minimum height=1.45cm] (P3) at (-3.0,-5.995) {};
    \node[texto, text width=2.5cm] at (-3.90,-5.995)
      {1.c~~Predecir $\mathrm{MDNBR}_{\min}$ con la DNN};
    \pic at (-2.12,-5.995) {dnnmini};

    %-------------------------------------------------- rama de techo
    \node[subcaja] (T1) at (3.0,-3.07)
      {2.a~~Establecer el flujo\\ de calor m\'aximo $q''_{\max}$};
    \node[subcaja] (T2) at (3.0,-4.42)
      {2.b~~Calcular el caudal\\ m\'asico $G$ para $q''_{\max}$};
    \node[subcaja, minimum height=1.45cm] (T3) at (3.0,-5.995) {};
    \node[texto, text width=2.5cm] at (2.10,-5.995)
      {2.c~~Predecir $\mathrm{MDNBR}_{\max}$ con la DNN};
    \pic at (3.88,-5.995) {dnnmini};

    %-------------------------------------------------- lazo de la secante
    \node[caja] (C) at (0,-8.40)
      {3.~M\'etodo de la secante:\\
       nuevo flujo de calor $q''_{k+1}$};

    \node[caja, minimum height=0.85cm] (D1) at (0,-9.675)
      {3.a~~Calcular el caudal m\'asico $G_{k+1}$};

    \node[caja, minimum height=1.25cm] (D2) at (0,-11.075) {};
    \node[texto, text width=3.6cm] at (-0.92,-11.075)
      {3.b~~Predecir $\mathrm{MDNBR}_{k+1}$ con la DNN};
    \pic at (1.45,-11.075) {dnnmini};

    %-------------------------------------------------- criterio de convergencia
    \draw[rombo] (0,-12.05) -- (4.2,-13.65) -- (0,-15.25) -- (-4.2,-13.65)
      -- cycle;
    \node[texto] at (0,-13.65)
      {$\begin{aligned}
         \bigl|\mathrm{MDNBR}_{k+1}-\mathrm{MDNBR}^{*}\bigr|
           &\le\varepsilon_{\mathrm{M}} \ \text{y} \\[-1pt]
         \bigl|G_{k+1}-G_{k}\bigr| &\le\varepsilon_{G} \ \text{y} \\[-1pt]
         \bigl|q''_{k+1}-q''_{k}\bigr| &\le\varepsilon_{q} \ ?
       \end{aligned}$};

    \node[final] (F) at (6.705,-13.65) {$q''_{\mathrm{PLC}}=q''_{k+1}$};

    %-------------------------------------------------- flujo: 0 -> ramas
    \draw[linea]  (A.south) -- (0,-0.95);
    \draw[linea]  (-3.0,-0.95) -- (3.0,-0.95);
    \draw[flecha] (-3.0,-0.95) -- (-3.0,-1.35);
    \draw[flecha] ( 3.0,-0.95) -- ( 3.0,-1.35);

    %-------------------------------------------------- flujo dentro de cada rama
    \draw[flecha] (P1.south) -- (P2.north);
    \draw[flecha] (P2.south) -- (P3.north);
    \draw[flecha] (T1.south) -- (T2.north);
    \draw[flecha] (T2.south) -- (T3.north);

    %-------------------------------------------------- ramas -> secante
    \draw[linea]  (-3.0,-6.97) -- (-3.0,-7.50) -- (0,-7.50);
    \draw[linea]  ( 3.0,-6.97) -- ( 3.0,-7.50) -- (0,-7.50);
    \draw[flecha] (0,-7.50) -- (C.north);

    %-------------------------------------------------- lazo
    \draw[flecha] (C.south)  -- (D1.north);
    \draw[flecha] (D1.south) -- (D2.north);
    \draw[flecha] (D2.south) -- (0,-12.05);

    \draw[flecha] (4.2,-13.65) -- (F.west);
    \node[nombre] at (4.75,-13.20) {s\'i};

    \draw[linea]  (-4.2,-13.65) -- (-5.7,-13.65) -- (-5.7,-8.40);
    \draw[flecha] (-5.7,-8.40) -- (C.west);
    \node[nombre, anchor=west] at (-5.3,-11.0) {no};
  \end{tikzpicture}}
  \caption{Lazo de c\'alculo del flujo de calor que alcanza el MDNBR objetivo.}
  \label{fig:dnb-dnn-loop}
\end{figure}
.
% Requiere en el preambulo del documento principal:
%   \usepackage{tikz}
%   \usetikzlibrary{arrows.meta}
%   \usepackage{amsmath,amssymb}
%
% Estructura: condiciones de operacion -> ramas paralelas de piso (q''_min) y
% techo (q''_max), cada una con establecer flujo, calcular caudal masico y
% predecir MDNBR con la DNN -> lazo de la secante, con el calculo del caudal
% masico y la prediccion de la DNN como pasos separados, hasta cumplir el
% criterio de convergencia.
%
% Los tres pasos que invocan a la red llevan un esquema en miniatura de la DNN
% (pic "dnnmini"), consistente con fig:mlp-arquitectura.

\begin{figure}[htbp]
  \centering
  \resizebox{\textwidth}{!}{
  \begin{tikzpicture}[
      caja/.style   = {rectangle, rounded corners=2pt, draw=black!70,
                       fill=black!4, line width=0.6pt, text width=6.0cm,
                       align=center, font=\footnotesize, minimum height=1.0cm,
                       inner sep=3pt},
      subcaja/.style= {rectangle, rounded corners=2pt, draw=black!70,
                       fill=white, line width=0.6pt, text width=4.6cm,
                       align=center, font=\footnotesize, minimum height=1.0cm,
                       inner sep=3pt},
      final/.style  = {rectangle, rounded corners=2pt, draw=black!80,
                       fill=black!8, line width=0.9pt, text width=2.6cm,
                       align=center, font=\footnotesize, minimum height=0.85cm,
                       inner sep=3pt},
      grupo/.style  = {rectangle, rounded corners=3pt, draw=black!45,
                       line width=0.6pt, dashed},
      titulo/.style = {font=\footnotesize\bfseries, inner sep=1pt},
      rombo/.style  = {draw=black!70, fill=black!4, line width=0.6pt},
      texto/.style  = {font=\footnotesize, align=center, inner sep=1pt},
      nombre/.style = {font=\footnotesize, inner sep=1pt},
      flecha/.style = {-{Latex[length=2mm,width=1.6mm]}, draw=black!70,
                       line width=0.6pt},
      linea/.style  = {draw=black!70, line width=0.6pt},
      %--- estilos del esquema en miniatura de la red
      nnnodo/.style = {circle, draw=black!65, fill=white, line width=0.3pt,
                       minimum size=1.4mm, inner sep=0pt},
      nnarco/.style = {draw=black!30, line width=0.25pt},
      %--- esquema en miniatura de la DNN (3 entradas, 4+4 ocultas, 2 salidas)
      dnnmini/.pic  = {
        \foreach \ya in {0.20,0,-0.20}
          \foreach \yb in {0.30,0.10,-0.10,-0.30}
            \draw[nnarco] (0,\ya) -- (0.4,\yb);
        \foreach \ya in {0.30,0.10,-0.10,-0.30}
          \foreach \yb in {0.30,0.10,-0.10,-0.30}
            \draw[nnarco] (0.4,\ya) -- (0.8,\yb);
        \foreach \ya in {0.30,0.10,-0.10,-0.30}
          \foreach \yb in {0.10,-0.10}
            \draw[nnarco] (0.8,\ya) -- (1.2,\yb);
        \foreach \y in {0.20,0,-0.20}
          \node[nnnodo] at (0,\y) {};
        \foreach \y in {0.30,0.10,-0.10,-0.30} {
          \node[nnnodo] at (0.4,\y) {};
          \node[nnnodo] at (0.8,\y) {};
        }
        \foreach \y in {0.10,-0.10}
          \node[nnnodo] at (1.2,\y) {};
      },
    ]
    \useasboundingbox (-6.25,-15.80) rectangle (8.61,1.0);

    %-------------------------------------------------- 0. condiciones
    \node[caja] (A) at (0,0)
      {0.~Definir las condiciones de operaci\'on\\
       y el objetivo $\mathrm{MDNBR}^{*}$};

    %-------------------------------------------------- marcos de las ramas
    \draw[grupo] (-5.7,-6.97) rectangle (-0.3,-1.35);
    \draw[grupo] ( 0.3,-6.97) rectangle ( 5.7,-1.35);

    \node[titulo] at (-3.0,-1.96) {1.~C\'alculo de piso};
    \node[titulo] at ( 3.0,-1.96) {2.~C\'alculo de techo};

    %-------------------------------------------------- rama de piso
    \node[subcaja] (P1) at (-3.0,-3.07)
      {1.a~~Establecer el flujo\\ de calor m\'inimo $q''_{\min}$};
    \node[subcaja] (P2) at (-3.0,-4.42)
      {1.b~~Calcular el caudal\\ m\'asico $G$ para $q''_{\min}$};
    \node[subcaja, minimum height=1.45cm] (P3) at (-3.0,-5.995) {};
    \node[texto, text width=2.5cm] at (-3.90,-5.995)
      {1.c~~Predecir $\mathrm{MDNBR}_{\min}$ con la DNN};
    \pic at (-2.12,-5.995) {dnnmini};

    %-------------------------------------------------- rama de techo
    \node[subcaja] (T1) at (3.0,-3.07)
      {2.a~~Establecer el flujo\\ de calor m\'aximo $q''_{\max}$};
    \node[subcaja] (T2) at (3.0,-4.42)
      {2.b~~Calcular el caudal\\ m\'asico $G$ para $q''_{\max}$};
    \node[subcaja, minimum height=1.45cm] (T3) at (3.0,-5.995) {};
    \node[texto, text width=2.5cm] at (2.10,-5.995)
      {2.c~~Predecir $\mathrm{MDNBR}_{\max}$ con la DNN};
    \pic at (3.88,-5.995) {dnnmini};

    %-------------------------------------------------- lazo de la secante
    \node[caja] (C) at (0,-8.40)
      {3.~M\'etodo de la secante:\\
       nuevo flujo de calor $q''_{k+1}$};

    \node[caja, minimum height=0.85cm] (D1) at (0,-9.675)
      {3.a~~Calcular el caudal m\'asico $G_{k+1}$};

    \node[caja, minimum height=1.25cm] (D2) at (0,-11.075) {};
    \node[texto, text width=3.6cm] at (-0.92,-11.075)
      {3.b~~Predecir $\mathrm{MDNBR}_{k+1}$ con la DNN};
    \pic at (1.45,-11.075) {dnnmini};

    %-------------------------------------------------- criterio de convergencia
    \draw[rombo] (0,-12.05) -- (4.2,-13.65) -- (0,-15.25) -- (-4.2,-13.65)
      -- cycle;
    \node[texto] at (0,-13.65)
      {$\begin{aligned}
         \bigl|\mathrm{MDNBR}_{k+1}-\mathrm{MDNBR}^{*}\bigr|
           &\le\varepsilon_{\mathrm{M}} \ \text{y} \\[-1pt]
         \bigl|G_{k+1}-G_{k}\bigr| &\le\varepsilon_{G} \ \text{y} \\[-1pt]
         \bigl|q''_{k+1}-q''_{k}\bigr| &\le\varepsilon_{q} \ ?
       \end{aligned}$};

    \node[final] (F) at (6.705,-13.65) {$q''_{\mathrm{PLC}}=q''_{k+1}$};

    %-------------------------------------------------- flujo: 0 -> ramas
    \draw[linea]  (A.south) -- (0,-0.95);
    \draw[linea]  (-3.0,-0.95) -- (3.0,-0.95);
    \draw[flecha] (-3.0,-0.95) -- (-3.0,-1.35);
    \draw[flecha] ( 3.0,-0.95) -- ( 3.0,-1.35);

    %-------------------------------------------------- flujo dentro de cada rama
    \draw[flecha] (P1.south) -- (P2.north);
    \draw[flecha] (P2.south) -- (P3.north);
    \draw[flecha] (T1.south) -- (T2.north);
    \draw[flecha] (T2.south) -- (T3.north);

    %-------------------------------------------------- ramas -> secante
    \draw[linea]  (-3.0,-6.97) -- (-3.0,-7.50) -- (0,-7.50);
    \draw[linea]  ( 3.0,-6.97) -- ( 3.0,-7.50) -- (0,-7.50);
    \draw[flecha] (0,-7.50) -- (C.north);

    %-------------------------------------------------- lazo
    \draw[flecha] (C.south)  -- (D1.north);
    \draw[flecha] (D1.south) -- (D2.north);
    \draw[flecha] (D2.south) -- (0,-12.05);

    \draw[flecha] (4.2,-13.65) -- (F.west);
    \node[nombre] at (4.75,-13.20) {s\'i};

    \draw[linea]  (-4.2,-13.65) -- (-5.7,-13.65) -- (-5.7,-8.40);
    \draw[flecha] (-5.7,-8.40) -- (C.west);
    \node[nombre, anchor=west] at (-5.3,-11.0) {no};
  \end{tikzpicture}}
  \caption{Lazo de c\'alculo del flujo de calor que alcanza el MDNBR objetivo.}
  \label{fig:dnb-dnn-loop}
\end{figure}
